\documentclass[a4paper, 11pt]{article}

\usepackage[T1,T2A]{fontenc}
\usepackage[utf8]{inputenc}
\usepackage[english, russian]{babel}

\usepackage{comment}
\usepackage{fullpage}
\usepackage{amssymb}
\usepackage{amsmath}
\usepackage{amsthm}
\usepackage{enumitem}

\begin{document}

\begin{center} \Large\bfseries Листок 11\\ Тема 11 (4.1). $p$-адические числа: элементарное определение и свойства \end{center}

\noindent\textbf{Упражнения и задачи}

\begin{enumerate}[topsep=0pt]

    \item Докажите, что различные канонические последовательности определяют различные целые $p$-адические числа. %[БШ]

    \item Докажите, что определённые на лекции сложение и умножение в $\mathbb{Z}_p$ не зависят от выбора последовательностей-представителей $\alpha\sim(x_n)$, $\beta\sim(y_n)$. Проверьте, что $\mathbb{Z}_p$ является коммутативным кольцом. %[БШ]

    \item Пусть $\alpha = \sum_{n=0}^{\infty} a_n p^n \in \mathbb{Z}_p$, где $0\le a_n<p$. Какой вид имеет разложение числа $-\alpha$? %[БШ, G]

    \item Докажите, что для любых $\alpha\in\mathbb{Z}_p$ и $n\ge 1$ существует $a\in\mathbb{Z}$, для которого $\alpha\equiv a\pmod{p^n}$. Докажите также, что для $a,b\in\mathbb{Z}$ сравнение $a\equiv b\pmod{p^n}$ в $\mathbb{Z}_p$ равносильно тому же сравнению в $\mathbb{Z}$. %[БШ]

    \item Докажите, что уравнение $x^2=2$ не имеет решений в $\mathbb{Q}_5$.

    \item Пусть $p\neq 2$, $m\in\mathbb{Z}$ и $(m,p)=1$. Сформулируйте и докажите необходимое и достаточное условие разрешимости уравнения $x^2=m$ в $\mathbb{Q}_p$. Сделайте вывод, что $\mathbb{Q}_p$ не является алгебраически замкнутым.

    \item Докажите, что если $\xi_n\to\xi\neq 0$ в $\mathbb{Q}_p$, то $\xi_n\neq 0$ при всех достаточно больших $n$ и $1/\xi_n\to 1/\xi$. %[БШ]

    \item Докажите $p$-адический аналог утверждения из анализа: из всякой ограниченной в $p$-адической норме последовательности в $\mathbb{Q}_p$ можно выделить сходящуюся подпоследовательность. %[БШ]

    \item Докажите $p$-адический критерий Коши: последовательность $(\xi_n)$ в $\mathbb{Q}_p$ сходится тогда и только тогда, когда для любого $N\ge 1$ существует $n_0$ такое, что $\nu_p(\xi_m-\xi_n)\ge N$ при всех $m,n\ge n_0$. %[БШ]

    \item Пусть $c \in \mathbb{Z}$, $p \nmid c$. Докажите, что последовательность $(c^{p^n})_{n\ge0}$ сходится в $\mathbb{Z}_p$, при этом для $\gamma = \lim c^{p^n}$ имеем $\gamma \equiv c \pmod{p}$, $\gamma^{p-1}=1$. Используя этот результат, докажите, что многочлен $t^{p-1}-1$ раскладывается на линейные множители в $\mathbb Z_p[t]$.

    \item Пусть последовательность $(x_n)$ определена как $x_n = 1+p+\dots +p^{n-1}$. Докажите, что $p$-адически  $x_n \rightarrow 1/(1-p)$, $n \rightarrow \infty$. %[БШ]
    
    \item Докажите, что для всякого $\xi\in\mathbb{Z}_p\cap\mathbb{Q}$ коэффициенты его разложения $\xi=\sum_{n=0}^{\infty}a_np^n$, $0\le a_n<p$, начиная с некоторого номера периодичны: существуют $k_0\ge 0$ и $m\ge 1$ такие, что $a_{k+m}=a_k$ при всех $k\ge k_0$. Докажите обратное: если коэффициенты такого разложения начиная с некоторого номера периодичны, то $\xi\in\mathbb{Q}$.
    
\end{enumerate}

% \noindent\textbf{SageMath}

% \begin{itemize}[topsep=0pt]
%     \item Исследуйте основные функции SageMath связанные с арифметикой $p$-адических чисел. Определение кольца и поле $p$-адических чисел: \texttt{Zp(n)}, \texttt{Qp(n)}. Рассмотрите примеры уравнения $x^2=m$ (используйте функцию \texttt{sqrt()}).
        
% \end{itemize}

\noindent\textbf{Темы для самостоятельного изучения}
\begin{itemize}[topsep=0pt]
    \item Элементы $p$-адического анализа: последовательности, ряды, дифференцирование, интегрирование. [Gouv], \S\S 5.1--5.4.
\end{itemize}

\end{document}